\section{Method}
\label{sec:method}

\subsection{Overview}
\label{subsec:overview}
Given a target specified in the first frame, a multimodal tracker localizes it in every subsequent frame using an RGB stream paired with one auxiliary modality $\mathcal{M}\!\in\!\{\text{Thermal},\text{Depth},\text{Event}\}$. 
%While standard trackers assume these streams are always complete, in practice the auxiliary stream is frequently \emph{incomplete}: sensors de-synchronize, drop frames, or fail entirely. 
While standard multimodal trackers typically assume complete and synchronized input streams, either stream, most often the auxiliary one, may become temporarily or persistently unavailable in practice due to sensor synchronization issues, self-calibration, or data transmission errors~\cite{lu2025modality,tan2025you}.
The conference version of this work, FlexTrack~\cite{tan2025you}, addressed this setting with a heterogeneous mixture-of-experts (HMoE) fusion, whose experts of heterogeneous capacity allow the router to spend more computation on harder, degraded inputs, and exposed the tracker to modality dropping via video-level random masking during training. However, it tolerates missing modalities passively: once a stream drops, its features are masked out.

In this paper, FlexTrack-V2 upgrades this design to active compensation through three coupled components. Bidirectional Modality Reconstruction (BMR) hallucinates the features of an absent stream from the present one in either direction, so that the inherited HMoE always routes over a completed feature pair.
Curriculum Missing Augmentation (CMA) replaces video-level, fixed-rate masking with per-sample, feature-level missing simulation whose severity grows over training.
Finally, a dual-stream self-distillation objective runs a modality-complete teacher fusion stream alongside the CMA-degraded student stream on a shared encoding, pulling degraded representations toward their complete counterparts.
Together they form a closed training loop: CMA synthesizes increasingly severe missing conditions, BMR fills the resulting feature gaps, and self-distillation supervises the compensation.

% In this paper, FlexTrack-V2 advances this adaptive-capacity principle by replacing the \emph{passive masking} of an absent modality with \emph{active reconstruction}, and stabilizes the training of this generative pathway through a completeness curriculum and online self-distillation.

The remainder of this section is organized as follows. 
Section~\ref{subsec:arch} introduces the modality-agnostic base architecture: a weight-shared Fast-iTPN encoder that tokenizes both streams identically, and a Mamba interaction neck that carries temporal state across frames.
Section~\ref{subsec:hmoe} reviews the HMoE fusion inherited from the conference version.
Sections~\ref{subsec:bmr} and~\ref{subsec:cma} detail the newly proposed BMR and CMA with dual-stream self-distillation, respectively.
Section~\ref{subsec:loss} summarizes the overall training objective.

% The proposed framework comprises four main components, as illustrated in Fig.~\ref{fig:framework}. First, Section~\ref{subsec:arch} introduces the unified tracking architecture, built upon a shared Fast-iTPN encoder and a Mamba interaction neck, which tokenizes and models long-range dependencies for any modality pair. Second, Section~\ref{subsec:hmoe} details the Heterogeneous Mixture-of-Experts (HMoE) fusion module, designed to dynamically allocate computation capacity based on modality availability. Third, Section~\ref{subsec:bmr} presents the Bilateral Modality-specific feature Reconstruction and Hallucination (BMR) module, which actively synthesizes missing-modality features from the available stream. Finally, Section~\ref{subsec:cma} proposes the Curriculum Missing Augmentation (CMA) with online self-distillation to align representations across different completeness levels. The overall training objective is summarized in Section~\ref{subsec:loss}.

\paragraph*{Notation}
For a search region, we denote the per-modality token features produced by the encoder as $\mathbf{f}^{r}\!\in\!\mathbb{R}^{N\times d}$ (RGB) and $\mathbf{f}^{a}\!\in\!\mathbb{R}^{N\times d}$ (auxiliary), where $N$ is the number of tokens and $d$ is the channel width. $\operatorname{sg}[\cdot]$ denotes the stop-gradient operator, and $t\!\in\!\{1,\dots,T\}$ indexes the training epoch.

\subsection{Unified Tracking Architecture}
\label{subsec:arch}

\paragraph*{Tokenization and encoder}
Multimodal tracking typically requires separate specialized encoders, which drastically increases parameter count. To achieve a unified framework, each input frame is a $6$-channel tensor formed by concatenating the RGB image with the (colorized) auxiliary map along the channel dimension. It is then split back into $\mathbf{x}^{r}$ and $\mathbf{x}^{a}$ before encoding. We adopt a Fast-iTPN backbone~\cite{tian2023fastitpn} as the encoder $E$: a three-stage pyramid where the first two stages are attention-free convolutional blocks and the third stage is a stack of self-attention blocks operating at stride $16$. The template set ($M\!=\!5$ exemplars at $112\!\times\!112$) and the search region ($224\!\times\!224$) are patch-embedded and tagged with learnable \emph{token-type} embeddings to distinguish template-foreground, template-background, and search tokens. The encoder is applied with \emph{shared weights} to both modalities:
\begin{equation}
  \mathbf{f}^{r}=E(\mathbf{x}^{r}),\qquad
  \mathbf{f}^{a}=E(\mathbf{x}^{a}),
  \label{eq:encode}
\end{equation}
This design keeps the parameter count independent of the auxiliary sensor type and allows a single model to serve RGB-T, RGB-D, and RGB-E tracking seamlessly.

\paragraph*{Interaction neck and head}
Standard self-attention mechanisms struggle to efficiently model long sequence dependencies required for temporal tracking. To address this, the fused search tokens (from Sec.~\ref{subsec:hmoe}) are refined by a Mamba interaction neck: a stack of $L\!=\!4$ selective state-space blocks~\cite{gu2023mamba} interleaved with injector/extractor cross-attention. This design acts similarly to a ViT-Adapter, efficiently feeding refined evidence back into the corresponding groups of encoder self-attention blocks with linear complexity. Finally, a center-based head~\cite{ye2022ostrack} predicts a classification score map, a local offset, and a size map, decoding the target box at the score-map peak.

\subsection{Heterogeneous Mixture-of-Experts Fusion}
\label{subsec:hmoe}
When fusing multimodal streams, complete and incomplete inputs demand different amounts of cross-modal computation. When both modalities are informative, fusion should exploit them fully. Conversely, when one is missing or degraded, an overweight fusion wastes capacity and injects noise. Previous MoE trackers use uniformly sized experts, failing to accommodate this dynamic requirement. We therefore propose fusing modalities with a \emph{heterogeneous} MoE whose experts have varying capacities. This enables the router to allocate more or fewer parameters per token according to the input quality.

Given the concatenated fusion input $\mathbf{h}=[\mathbf{f}^{r};\mathbf{f}^{a}]$, we use $K\!=\!8$ expert MLPs $\{E_i\}_{i=1}^{K}$ with exponentially growing hidden widths:
\begin{equation}
\begin{aligned}
  c_i &= 2^{\,i+1},\quad i=1,\dots,8, \\
      &\in \{4,8,16,32,64,128,256,512\}.
\end{aligned}
\label{eq:capacity}
\end{equation}
A lightweight expert cheaply processes easy or degraded tokens, while a heavy expert models rich cross-modal interactions. Routing follows noisy top-$k$ gating~\cite{shazeer2017moe} with $k\!=\!2$:
\begin{align}
  \mathbf{s} &= \mathbf{h}\mathbf{W}_g + \epsilon\odot\operatorname{softplus}(\mathbf{h}\mathbf{W}_n), \quad \epsilon\sim\mathcal{N}(0,\mathbf{I}),\\
  \mathbf{g} &= \operatorname{softmax}\!\big(\operatorname{TopK}(\mathbf{s}, 2)\big),
\end{align}
where $\mathbf{W}_g$ and $\mathbf{W}_n$ are learnable gating weights. The MoE routing loss $\mathcal{L}_{\text{moe}}$ encourages load balancing across experts and orthogonal gate directions:
\begin{equation}
  \mathcal{L}_{\text{moe}} = \operatorname{CV}\!\big(\text{imp}\big)^2 + \operatorname{CV}\!\big(\text{load}\big)^2 + \lambda_{\perp}\,\big\|\mathbf{W}_g^{\!\top}\mathbf{W}_g - \mathbf{I}\big\|_F^2,
  \label{eq:moe-loss}
\end{equation}
where $\operatorname{CV}(\cdot)$ is the coefficient of variation over the batch, $\text{imp}$ and $\text{load}$ are the per-expert routing importance and assignment counts, and the last term pushes the gate directions apart.

\subsection{Bilateral Modality Reconstruction and Hallucination}
\label{subsec:bmr}
Prior works typically mask absent modalities with zeros or duplicate the last available frame. This \emph{passive masking} discards implicit information: because multimodal streams of the same scene are highly correlated, the present modality inherently carries strong priors about the missing one. To actively compensate for missing data, we design a Bilateral Modality-specific feature Reconstruction (BMR) module. 

We introduce two lightweight cross-modal \emph{hallucinators}, $H_{r\rightarrow a}$ and $H_{a\rightarrow r}$, formulated as two-layer MLPs. Their output layers are zero-initialized to ensure training starts from an identity-preserving state. The hallucination process is defined as:
\begin{equation}
  \hat{\mathbf{f}}^{a}=H_{r\rightarrow a}(\mathbf{f}^{r}),\qquad
  \hat{\mathbf{f}}^{r}=H_{a\rightarrow r}(\mathbf{f}^{a}).
  \label{eq:hallucinate}
\end{equation}
Let $m^{r},m^{a}\!\in\!\{0,1\}$ represent the presence indicator for each modality in a given sample. Before routing into the fusion module, we substitute any missing stream with its hallucinated estimate:
\begin{equation}
  \tilde{\mathbf{f}}^{a}=m^{a}\,\mathbf{f}^{a}+(1-m^{a})\,\hat{\mathbf{f}}^{a},
  \label{eq:substitute}
\end{equation}
and symmetrically for $\tilde{\mathbf{f}}^{r}$. These substituted features feed directly into the HMoE described in Sec.~\ref{subsec:hmoe}, forming our joint \textbf{BMR-HMoE} pipeline. The hallucinators are supervised using an $L_2$ reconstruction loss only where the target modality is genuinely observed. A stop-gradient operator is applied to the target to prevent the reconstruction task from corrupting the shared encoder:
\begin{equation}
  \mathcal{L}_{\text{rec}} = m^{a}\big\|\hat{\mathbf{f}}^{a}-\operatorname{sg}[\mathbf{f}^{a}]\big\|_2^2 + m^{r}\big\|\hat{\mathbf{f}}^{r}-\operatorname{sg}[\mathbf{f}^{r}]\big\|_2^2.
  \label{eq:rec-loss}
\end{equation}
Compared to masking, BMR ensures the fusion network always receives complete semantic features, significantly preventing tracking collapse during long-term sensor failures.

\subsection{Curriculum Missing Augmentation and Self-Distillation}
\label{subsec:cma}

\begin{algorithm}[tb]
\caption{Curriculum Missing Augmentation (CMA) with Active Compensation}
\label{alg:cma}
\KwIn{Epoch $t$, Max epochs $T$, Batch size $B$, RGB tokens $\mathbf{f}^{r}$, Auxiliary tokens $\mathbf{f}^{a}$, Hallucinators $H$, HMoE module.}
\KwOut{Student fused features $\mathbf{F}_{\text{S}}$, Teacher fused features $\mathbf{F}_{\text{T}}$.}
\tcp{1. Compute dynamic missing probability (Eq.~\ref{eq:cma})}
$p_{\text{miss}}(t) \leftarrow p_{\min} + \min\!\Big(1,\frac{t-1}{T-1}\Big) (p_{\max} - p_{\min})$\;
\For{$b = 1$ \KwTo $B$}{
    \tcp{2. Generate Teacher (Complete) Features}
    $\mathbf{h}_{\text{T}} \leftarrow [\mathbf{f}^{r}_b; \mathbf{f}^{a}_b]$\;
    $\mathbf{F}_{\text{T},b} \leftarrow \text{HMoE}(\mathbf{h}_{\text{T}})$\;
    
    \tcp{3. Determine Curriculum Masking}
    $p \sim \mathcal{U}(0,1)$\;
    \eIf{$p < p_{\text{miss}}(t)$}{
        $M_b \leftarrow \operatorname{random}([[1,0], [0,1]])$ \tcp*{Randomly drop one modality}
        $m^{r}_b \leftarrow M_b[0], \quad m^{a}_b \leftarrow M_b[1]$\;
    }{
        $m^{r}_b \leftarrow 1, \quad m^{a}_b \leftarrow 1$ \tcp*{Keep both modalities intact}
    }
    
    \tcp{4. Bilateral Feature Hallucination (Eq.~\ref{eq:hallucinate})}
    $\hat{\mathbf{f}}^{a}_b \leftarrow H_{r\rightarrow a}(\mathbf{f}^{r}_b)$\;
    $\hat{\mathbf{f}}^{r}_b \leftarrow H_{a\rightarrow r}(\mathbf{f}^{a}_b)$\;
    
    \tcp{5. Active Substitution (Eq.~\ref{eq:substitute})}
    $\tilde{\mathbf{f}}^{a}_b \leftarrow m^{a}_b \mathbf{f}^{a}_b + (1-m^{a}_b) \hat{\mathbf{f}}^{a}_b$\;
    $\tilde{\mathbf{f}}^{r}_b \leftarrow m^{r}_b \mathbf{f}^{r}_b + (1-m^{r}_b) \hat{\mathbf{f}}^{r}_b$\;
    
    \tcp{6. Generate Student (Degraded) Features}
    $\mathbf{h}_{\text{S}} \leftarrow [\tilde{\mathbf{f}}^{r}_b; \tilde{\mathbf{f}}^{a}_b]$\;
    $\mathbf{F}_{\text{S},b} \leftarrow \text{HMoE}(\mathbf{h}_{\text{S}})$\;
}
\Return $\mathbf{F}_{\text{S}}, \mathbf{F}_{\text{T}}$\;
\end{algorithm}

Training a reconstruction model typically requires a rigid, frozen external teacher model. Furthermore, training only on complete data leaves the missing-regime pathways untrained, while immediately training with high missing rates destabilizes early optimization. To circumvent these issues, we propose Curriculum Missing Augmentation (CMA) with online self-distillation.

For each training sample, we randomly drop an auxiliary modality's search tokens, replacing them with a learned \emph{missing prompt}, using a probability $p_{\text{miss}}(t)$ that dynamically anneals over training epochs $t$:
\begin{equation}
  p_{\text{miss}}(t)=p_{\min}+\min\!\Big(1,\tfrac{t-1}{T-1}\Big)\,(p_{\max}-p_{\min}).
  \label{eq:cma}
\end{equation}
We linearly ramp the probability from $p_{\min}\!=\!0$ to $p_{\max}\!=\!0.35$ over $T\!=\!40$ epochs. To ensure the reconstructed representations align with complete-modality ones, we execute two forward passes of the \emph{same} model. A teacher pass processes the intact modality pair to yield fused features $\mathbf{F}_{\text{T}}$, while a student pass processes the curriculum-dropped pair to yield $\mathbf{F}_{\text{S}}$. We then apply a self-distillation loss:
\begin{equation}
  \mathcal{L}_{\text{dist}} = \big\|\mathbf{F}_{\text{S}}-\operatorname{sg}[\mathbf{F}_{\text{T}}]\big\|_2^2.
  \label{eq:distill}
\end{equation}
This formulation forces the model to act as its own online teacher. By binding the BMR hallucination targets to high-quality complete representations, CMA guarantees robust optimization without the heavy memory overhead of external teachers.

\subsection{Training Objective}
\label{subsec:loss}
The tracking supervision combines a focal classification loss $\mathcal{L}_{\text{cls}}$ on the score map with generalized-IoU and $\ell_1$ box-regression losses. The total loss seamlessly integrates the task supervision with the fusion, reconstruction, and distillation terms:
\begin{equation}
\begin{aligned}
  \mathcal{L} =\;& \lambda_{\text{giou}}\mathcal{L}_{\text{giou}} + \lambda_{\ell_1}\mathcal{L}_{\ell_1} + \lambda_{\text{cls}}\mathcal{L}_{\text{cls}} \\
  &+ \lambda_{\text{dist}}\mathcal{L}_{\text{dist}} + \lambda_{\text{fus}}\big(\mathcal{L}_{\text{moe}}+\mathcal{L}_{\text{rec}}\big).
\end{aligned}
\label{eq:total-loss}
\end{equation}
Unless otherwise specified, we empirically set $\lambda_{\text{giou}}\!=\!2$, $\lambda_{\ell_1}\!=\!5$, $\lambda_{\text{cls}}\!=\!2$, $\lambda_{\text{dist}}\!=\!2.5$, and $\lambda_{\text{fus}}\!=\!10$. The model is trained end-to-end using AdamW (learning rate $3\!\times\!10^{-4}$, weight decay $10^{-4}$, encoder learning rate scaled by $0.1$) for $T\!=\!40$ epochs with a $10\times$ step decay at epoch $25$. We use a batch size of $32$ and gradient-norm clipping at $0.1$, jointly trained over VisEvent, LasHeR, and DepthTrack datasets with equal sampling.
